% Options for packages loaded elsewhere
\PassOptionsToPackage{unicode}{hyperref}
\PassOptionsToPackage{hyphens}{url}
%
\documentclass[
  ignorenonframetext, 10pt
]{beamer}
\usepackage{pgfpages}
\setbeamertemplate{caption}[numbered]
\setbeamertemplate{caption label separator}{: }
\setbeamercolor{caption name}{fg=normal text.fg}
\beamertemplatenavigationsymbolsempty
% Prevent slide breaks in the middle of a paragraph
\widowpenalties 1 10000
\raggedbottom
\setbeamertemplate{part page}{
  \centering
  \begin{beamercolorbox}[sep=16pt,center]{part title}
    \usebeamerfont{part title}\insertpart\par
  \end{beamercolorbox}
}
\setbeamertemplate{section page}{
  \centering
  \begin{beamercolorbox}[sep=12pt,center]{part title}
    \usebeamerfont{section title}\insertsection\par
  \end{beamercolorbox}
}
\setbeamertemplate{subsection page}{
  \centering
  \begin{beamercolorbox}[sep=8pt,center]{part title}
    \usebeamerfont{subsection title}\insertsubsection\par
  \end{beamercolorbox}
}
\AtBeginPart{
  \frame{\partpage}
}
\AtBeginSection{
  \ifbibliography
  \else
    \frame{\sectionpage}
  \fi
}
\AtBeginSubsection{
  \frame{\subsectionpage}
}
\usepackage{amsmath,amssymb}
\usepackage{iftex}
\ifPDFTeX
  \usepackage[T1]{fontenc}
  \usepackage[utf8]{inputenc}
  \usepackage{textcomp} % provide euro and other symbols
\else % if luatex or xetex
  \usepackage{unicode-math} % this also loads fontspec
  \defaultfontfeatures{Scale=MatchLowercase}
  \defaultfontfeatures[\rmfamily]{Ligatures=TeX,Scale=1}
\fi
\usepackage{lmodern}
\ifPDFTeX\else
  % xetex/luatex font selection
\fi
% Use upquote if available, for straight quotes in verbatim environments
\IfFileExists{upquote.sty}{\usepackage{upquote}}{}
\IfFileExists{microtype.sty}{% use microtype if available
  \usepackage[]{microtype}
  \UseMicrotypeSet[protrusion]{basicmath} % disable protrusion for tt fonts
}{}
\makeatletter
\@ifundefined{KOMAClassName}{% if non-KOMA class
  \IfFileExists{parskip.sty}{%
    \usepackage{parskip}
  }{% else
    \setlength{\parindent}{0pt}
    \setlength{\parskip}{6pt plus 2pt minus 1pt}}
}{% if KOMA class
  \KOMAoptions{parskip=half}}
\makeatother
\usepackage{xcolor}
\newif\ifbibliography
\setlength{\emergencystretch}{3em} % prevent overfull lines
\providecommand{\tightlist}{%
  \setlength{\itemsep}{0pt}\setlength{\parskip}{0pt}}
\setcounter{secnumdepth}{-\maxdimen} % remove section numbering
%\usetheme{metropolis}


%\setbeamercolor{normal text}{bg=white}
%\addtobeamertemplate{frametitle}{}{\vspace{-0.5em}} % decrease
%\setbeamersize{text margin left=5mm,text margin right=10mm}
%\usefonttheme[onlymath]{serif}
%\usefonttheme{professionalfonts}
%%%%%%%%%%%%%%%%
% Packages
%%%%%%%%%%%%%%%%

\usepackage{geometry}
\usepackage{graphicx}
\usepackage{amssymb}
\usepackage{amsopn}
\usepackage{cancel}
\usepackage{epstopdf}
\usepackage{amsmath}  	% this permits text in eqnarray among other benefits
\usepackage{url}		% produces hyperlinks
\usepackage{hyperref}	% allows for color usage in tables
\usepackage[english]{babel}
%\usepackage[latin1]{inputenc}
\usepackage{colortbl}	% allows for color usage in tables
\usepackage{multirow}	% allows for rows that span multiple rows in tables
\usepackage{color}		% this package has a variety of color options
\usepackage{xcolor}
\usepackage{pgf}
\usepackage{calc}
\usepackage{ulem}
\usepackage{multicol}
\usepackage{booktabs}
\usepackage{textcomp}
%\usepackage[misc]{ifsym} % for the dice symbol \Cube{}
%\usepackage{txfonts}
%\usepackage{listings}
%\usepackage{tikz}
\usepackage{array}
\usepackage{wasysym}
\usepackage{fancyvrb}

%% change fontsize of R code
\ifdefined\Shaded
\let\oldShaded\Shaded
\let\endoldShaded\endShaded
\renewenvironment{Shaded}{\footnotesize\oldShaded}{\endoldShaded}
\fi
%% change fontsize of output
\let\oldverbatim\verbatim
\let\endoldverbatim\endverbatim
\renewenvironment{verbatim}{\footnotesize\oldverbatim}{\endoldverbatim}

%%%%%%%%%%%%%%%%
% Remove navigation symbols
%%%%%%%%%%%%%%%%

\setbeamertemplate{navigation symbols}{}

\setbeamertemplate{footline}[frame number]

%%%%%%%%%%%%%%%%
% User defined colors
%%%%%%%%%%%%%%%%

\xdefinecolor{oiB}{rgb}{0.22,0.52,0.72}
\definecolor{oiG}{rgb}{.298,.447,.114}
\xdefinecolor{hlblue}{rgb}{0.051,0.65,1}
\xdefinecolor{gray}{rgb}{0.5, 0.5, 0.5}
\xdefinecolor{darkGray}{rgb}{0.3, 0.3, 0.3}
\xdefinecolor{darkerGray}{rgb}{0.2, 0.2, 0.2}
\xdefinecolor{rubineRed}{rgb}{0.89,0,0.30}
\xdefinecolor{irishGreen}{rgb}{0,0.60,0}	
\definecolor{lightGreen}{rgb}{0.387,0.581,0.148}

%%%%%%%%%%%%%%%%
% Template colors
%%%%%%%%%%%%%%%%

%\setbeamercolor*{palette primary}{fg=white,bg= oiB!80!black!90}
%\setbeamercolor*{palette secondary}{fg=black,bg= oiB!80!black}
%\setbeamercolor*{palette tertiary}{fg=white,bg= oiB!80!black!80}
%\setbeamercolor*{palette quaternary}{fg=white,bg= oiB}
%\setbeamercolor{structure}{fg= oiB}
%\setbeamercolor{frametitle}{bg= oiB!90}
%\setbeamertemplate{blocks}[shadow=false]
%\setbeamersize{text margin left=2em,text margin right=2em}

%\setbeamercolor{code body}{bg=gray!20!white!80,fg=black}


%%%%%%%%%%%%%%%%
% Get rid of fancy enumerated list bullets
%%%%%%%%%%%%%%%%

%\setbeamertemplate{enumerate items}[default]

%%%%%%%%%%%%%%%%
% Custom commands
%%%%%%%%%%%%%%%%

\newcommand{\Sum}{\sum\nolimits}
\newcommand{\Hide}[2]{\underline{~~~\onslide<#1>{\alert{#2}}~~~}}
\newcommand{\hide}[2]{\underline{~~\onslide<#1>{\alert{#2}}~~}}
\newcommand{\hid}[2]{\onslide<#1>{\alert{#2}}}
\newcommand{\code}[1]{\textcolor{blue}{\texttt{#1}}}

% red
\newcommand{\red}[1]{\textit{\textcolor{rubineRed}{#1}}}

% pink
\newcommand{\pink}[1]{\textit{\textcolor{rubineRed!90!white!50}{#1}}}

% green
\newcommand{\green}[1]{\textit{\textcolor{irishGreen}{#1}}}

% orange
\newcommand{\orange}[1]{\textit{\textcolor{orange}{#1}}}

% links: webURL, webLin, appLink
\newcommand{\webURL}[1]{\urlstyle{same}{ \textit{\textcolor{darkGray}{\url{#1}}}}}
\newcommand{\webLink}[2]{\href{#1}{\textcolor{darkGray}{{#2}}}}
\newcommand{\appLink}[2]{\href{#1}{\textcolor{white}{{#2}}}}

% mail
\newcommand{\mail}[1]{\href{mailto:#1}{\textit{\textcolor{darkGray}{#1}}}}

% highlighting: hl, hlGr, mathhl
\newcommand{\hl}[1]{\textit{\textcolor{hlblue}{#1}}}
\newcommand{\hlGr}[1]{\textit{\textcolor{lightGreen}{#1}}}
\newcommand{\mathhl}[1]{\textcolor{hlblue}{\ensuremath{#1}}}

% two col: two columns
\newenvironment{twocol}[4]{
\begin{columns}[c]
\column{#1\textwidth}
#3
\column{#2\textwidth}
#4
\end{columns}
}





%%%%%%%%%%%%%%%%
% Change margin
%%%%%%%%%%%%%%%%

\newenvironment{changemargin}[2]{%
\begin{list}{}{%
\setlength{\topsep}{0pt}%
\setlength{\leftmargin}{#1}%
\setlength{\rightmargin}{#2}%
\setlength{\listparindent}{\parindent}%
\setlength{\itemindent}{\parindent}%
\setlength{\parsep}{\parskip}%
}%
\item}{\end{list}}

%%%%%%%%%%%%%%%%
% Footnote
%%%%%%%%%%%%%%%%

\long\def\symbolfootnote[#1]#2{\begingroup%
\def\thefootnote{\fnsymbol{footnote}}\footnote[#1]{#2}\endgroup}

%%%%%%%%%%%%%%%%
% Graphics
%%%%%%%%%%%%%%%%

\DeclareGraphicsRule{.tif}{png}{.png}{`convert #1 `dirname #1`/`basename #1 .tif`.png}



\setbeamercolor{normal text}{bg=white}

\DeclareMathOperator{\var}{var}
\newcommand{\boxmodel}[1]{\begin{array}{|c|}#1\\[2pt]\hline\end{array}}
\newcommand{\BXX}[2]{\begin{array}{|@{\,}c@{\,}|@{\,}c@{\,}|}\hline#1&#2\\[-1pt]\hline\end{array}}
\newcommand{\p}{\mathrm{P}}
\newcommand{\SE}{\mathrm{SE}}
\newcommand{\xbar}{\overline{x}}
\newcommand{\Xbar}{\overline{X}}
\newcommand{\ybar}{\overline{y}}
\newcommand{\Ybar}{\overline{Y}}
\newcommand{\zbar}{\overline{z}}
\newcommand{\Zbar}{\overline{Z}}
\newcommand{\wbar}{\overline{w}}
\newcommand{\Wbar}{\overline{W}}
\newcommand{\vbar}{\overline{v}}
\newcommand{\Vbar}{\overline{V}}
\newcommand{\dbar}{\overline{d}}
\newcommand{\Dbar}{\overline{D}}
\newcommand{\ebar}{\overline{e}}
\newcommand{\ehat}{\widehat{\epsilon}}
\newcommand{\yhat}{\widehat{y}}
\newcommand{\Yhat}{\widehat{Y}}
\newcommand{\bhat}{\widehat{\beta}}
\newcommand{\dl}[1]{\widehat{\delta}_{#1}}
\newcommand{\g}[1]{\widehat{\gamma}_{#1}}
\newcommand{\al}[1]{\widehat{\alpha}_{#1}}
\newcommand{\Bj}{{\beta_j}}
\newcommand{\Bp}{{\beta_p}}
\newcommand{\Bq}{{\beta_q}}
\renewcommand{\b}[1]{\widehat{\beta}_{#1}}
\newcommand{\B}[1]{{\beta}_{#1}}
\newcommand{\BETA}{{\boldsymbol\beta}}
\newcommand{\bp}{\widehat{\beta}_p}
\newcommand{\bq}{\widehat{\beta}_q}
\newcommand{\bj}{\widehat{\beta}_j}
\newcommand{\df}{\mathrm{df}}
\DeclareMathOperator{\VIF}{VIF}
\DeclareMathOperator{\SSE}{SSE}
\DeclareMathOperator{\SST}{SST}
\DeclareMathOperator{\SSR}{SSR}
\DeclareMathOperator{\MSE}{MSE}
\DeclareMathOperator{\MST}{MST}
\DeclareMathOperator{\MSR}{MSR}
\DeclareMathOperator{\RM}{RM}
\DeclareMathOperator{\FM}{FM}
\DeclareMathOperator{\E}{E}
\DeclareMathOperator{\SD}{SD}
\DeclareMathOperator{\se}{se}
\DeclareMathOperator{\cov}{Cov}
\DeclareMathOperator{\corr}{Corr}
\DeclareMathOperator{\cor}{Cor}
\DeclareMathOperator{\Var}{Var}
\DeclareMathOperator{\logit}{logit}
\newcommand{\hsigma}{\widehat{\sigma}}
\newcommand{\betahat}{\widehat{\beta}}
\newcommand{\slide}{<br><br><br><br><br>}
\newcommand{\bbeta}{\boldsymbol{\beta}}
\newcommand{\bX}{\mathbf{X}}
\newcommand{\bY}{\mathbf{Y}}
\newcommand{\bH}{\mathbf{H}}
\newcommand{\bI}{\mathbf{I}}
\newcommand{\be}{\mathbf{e}}
\newcommand{\pval}{\mathit{P}\mathrm{-val}}



\newcommand{\vmu}{\mbox{\boldmath{$\mu$}}}
\newcommand{\vR}{\mbox{\boldmath{$R$}}}
\newcommand{\vpsi}{\mbox{\boldmath{$\psi$}}}
\newcommand{\vZ}{\mbox{\boldmath{$Z$}}}
\newcommand{\vX}{\mbox{\boldmath{$X$}}}
\newcommand{\vY}{\mbox{\boldmath{$Y$}}}
\newcommand{\vW}{\mbox{\boldmath{$W$}}}
\newcommand{\vz}{\mbox{\boldmath{$z$}}}
\newcommand{\vr}{\mbox{\boldmath{$r$}}}
\newcommand{\vx}{\mbox{\boldmath{$x$}}}
\newcommand{\vk}{\mbox{\boldmath{$k$}}}
\newcommand{\vp}{\mbox{\boldmath{$p$}}}
\newcommand{\vv}{\mbox{\boldmath{$v$}}}
\newcommand{\vg}{\mbox{\boldmath{$g$}}}
\newcommand{\vbeta}{\mbox{\boldmath{$\beta$}}}
\newcommand{\vone}{\mbox{\boldmath{$1$}}}
\newcommand{\vep}{\mbox{\boldmath{$\epsilon$}}}
\newcommand{\veta}{\mbox{\boldmath{$\eta$}}}
\newcommand{\vl}{\mbox{\boldmath{$l$}}}
\newcommand{\ve}{\mbox{\boldmath{$e$}}}
\newcommand{\va}{\mbox{\boldmath{$a$}}}
\newcommand{\vb}{\mbox{\boldmath{$b$}}}
\newcommand{\vPsi}{\mbox{\boldmath{$\Psi$}}}
\DeclareMathOperator*{\V}{\mathbb{V}}


% geometry: topsep=0pt 
\ifLuaTeX
  \usepackage{selnolig}  % disable illegal ligatures
\fi
\IfFileExists{bookmark.sty}{\usepackage{bookmark}}{\usepackage{hyperref}}
\IfFileExists{xurl.sty}{\usepackage{xurl}}{} % add URL line breaks if available
\urlstyle{same}
\hypersetup{
  pdftitle={Chapter 1: Controlled Experiments},
  hidelinks,
  pdfcreator={LaTeX via pandoc}}

\title{Chapter 1: Controlled Experiments}
\author{}
\date{\vspace{-2.5em}}

\begin{document}
\frame{\titlepage}

\begin{frame}{}
\protect\hypertarget{section}{}
\underline{General Question}: How to test whether or not a new drug, for
example, a new vaccine, is effective?

\emph{Ask the class in small groups to come up with answers to this.}
\bigskip
\pause
Basic method: \textbf{Comparison}

\begin{itemize}
\tightlist
\item
  Divide people (subjects) into two groups;
\item
  Give the vaccine to one group (\red{treatment group}) and not to the
  other one (\red{control group});
\item
  Compare the responses (e.g.~the percentage of subjects getting the
  disease) in each group.

  \begin{itemize}
  \tightlist
  \item
    If the percentage is lower in the treatment group, then the vaccine
    is effective; otherwise not effective.
  \end{itemize}
\end{itemize}
\end{frame}

\begin{frame}{Example: Polio Vaccine (FPP)}
\protect\hypertarget{example-polio-vaccine-fpp}{}
\begin{tabular}{@{}ll}
 Polio & - contagious and mostly a child disease \\
 & - major epidemic hit U.S. in 1916 \\
 Vaccine & - developed by Salk (1950s)
 \end{tabular}
 \vspace{0.17in}

\underline{Question}: How to test its effectiveness? \vspace{0.07in}
\pause

\begin{itemize}
\tightlist
\item
  \textbf{Idea \#1:} vaccinate every child, see if polio incidence
  declines

  \begin{itemize}
  \tightlist
  \item
    \textbf{Issue:} Incidence fluctuates year to year.\pause
  \end{itemize}
\item
  \textbf{Idea \#2: } use comparison method and leave some children
  unvaccinated as a control\pause

  \begin{itemize}
  \tightlist
  \item
    \textbf{Issue:} Is this ethical?
  \item
    \textbf{Issue:} How to choose which children are vaccinated?\pause
  \end{itemize}
\item
  \textbf{Idea \#3:} use comparison method and vaccinate every child
  whose parents consent\pause

  \begin{itemize}
  \tightlist
  \item
    \textbf{Issue:} ?
  \end{itemize}
\end{itemize}
\end{frame}

\begin{frame}{Example: Polio Vaccine (First Experiment)}
\protect\hypertarget{example-polio-vaccine-first-experiment}{}
NFIP divided the subjects by school grade. \newline NFIP vaccinated all
grade 2 children whose parents consented.

\begin{itemize}
\tightlist
\item
  Control: Grades 1 \& 3;
\item
  Treatment: Grade 2
\end{itemize}

\begin{center}
\begin{tabular}{lcccc}
  &Grade & Size & Number   & Rate\\[-3pt]
  &      &      & Infected & \footnotesize (per 100,000)\\
\midrule
Treatment  &  2      & 225,000 & ~56 & 25\\
Control    &  1 \& 3 & 725,000 & 392 & 54\\
Not consent&  2      & 125,000 & ~55 & 44
\end{tabular}
\end{center}
\end{frame}

\begin{frame}{}
\protect\hypertarget{section-1}{}
Should we

\begin{itemize}
\tightlist
\item
  compare Treatment v.s Control?

  \begin{itemize}
  \tightlist
  \item
    why comparing rates, not counts?
  \item
    other issue?
  \end{itemize}
\item
  compare Treatment (Consent) v.s. (Not Consent)?
\end{itemize}
\end{frame}

\begin{frame}{Why Compare the \textbf{Rates}?}
\protect\hypertarget{why-compare-the}{}
In the U.S., the number of deaths from all causes is

\begin{itemize}
\tightlist
\item
  2.1 million in 1990;
\item
  1.7 million in 1960
\end{itemize}

nearly a 25\% increase.

\begin{block}{True or false, and why:}
The data show that the public's health got worse over
the period 1960-1990. (The population in 1960 was about 180 million,
in 1990 about 248 million)
\end{block}
\pause
\structure{False. The population got bigger too.\\
Death rate in 1960 $=1.7/180\times 100\%=0.94\%$\\
Death rate in 1990 $=2.1/248\times 100\%=0.85\%$\\
The death rate actually dropped over the period 1960-1990.
}
\end{frame}

\begin{frame}{}
\protect\hypertarget{section-2}{}
Why is it important that the treatment group and the control group are
similar? \bigskip\pause

\underline{Key to the comparison method}:

The two groups should be similar, and the only difference should be:
\textbf{treated or not} \bigskip\pause

Some \textbf{bad} ideas on how to divide subjects:

\begin{itemize}
\tightlist
\item
  by health, wealth, or gender;
\item
  by age, or school grade;
\item
  by volunteers/non-volunteers;
\end{itemize}

\bigskip

\pause

\begin{block}{Definition: Confounders/Confounding Variables}
Factors that can also explain the difference in the outcomes
of the treatment group and the control group
\end{block}
\end{frame}

\begin{frame}{Polio Vaccine: Volunteers/Non-Volunteers}
\protect\hypertarget{polio-vaccine-volunteersnon-volunteers}{}
\begin{itemize}
\tightlist
\item
  Need parents' permission to vaccinate their kids
\item
  If treatment = volunteers, control = non-volunteers,

  \begin{itemize}
  \tightlist
  \item
    Higher income parents would more likely consent to vaccinating their
    child;
  \item
    Higher income \(\rightarrow\) more hygenic environment
    \(\rightarrow\) less chance to contract mild cases of polio early in
    childhood \(\rightarrow\) less chance to develop antibodies to
    polio;
  \item
    So, volunteers tend to be more vulnerable to polio;
  \item
    The results will be \textbf{biased against} the new vaccine;
  \end{itemize}
\item
  The effect of being a volunteer/non-volunteer will be
  \textbf{confounded} with the effect of the treatment.
\item
  Conclusion: To avoid \textbf{confounding}, \red{both} the treatment
  group and control group must be chosen from volunteers;
\end{itemize}
\end{frame}

\begin{frame}{Possible Strategies to Deal with Confounders}
\protect\hypertarget{possible-strategies-to-deal-with-confounders}{}
For example, \textbf{age} is confounded with the effect of the vaccine:

\begin{itemize}
\tightlist
\item
  \textbf{restrict} the confounder (e.g.~study only 6-yr-old kids)
\item
  \textbf{control for} the confounder (e.g.~conduct separate experiments
  for each age group)
\item
  \textbf{balance} (e.g.~make the age compositions in the treatment and
  control group similar) \bigskip\pause
\end{itemize}

A simple way to ensure balance: \red{Randomization},\newline which means
assigning subjects to the two groups \red{randomly}.

\begin{itemize}
\tightlist
\item
  Randomly \(\neq\) Haphazardly
\item
  Avoid the human factor: use coin tossing/random number generation.
\end{itemize}
\end{frame}

\begin{frame}{Example: Polio Vaccine (Randomized Experiment)}
\protect\hypertarget{example-polio-vaccine-randomized-experiment}{}
\begin{itemize}
\tightlist
\item
  Of the children whose parents consented to vaccination, half were
  randomly assigned to the control group and the other half to the
  treatment group.

  \begin{itemize}
  \tightlist
  \item
    Hence this is a \textbf{randomized controlled} experiment.
  \end{itemize}
\item
  \textbf{Blinding}: The children in the control group were given a
  \textbf{placebo}.

  \begin{itemize}
  \tightlist
  \item
    Ensures response is to the vaccine, not the \red{idea} of
    vaccination.
  \end{itemize}
\item
  Doctors were unaware of which group each child was in.

  \begin{itemize}
  \tightlist
  \item
    This is known as \textbf{double blinding} since neither the child
    nor the doctor knows whether they got the vaccine or the placebo.
  \end{itemize}
\end{itemize}
\end{frame}

\begin{frame}{Blinding \& Placebo}
\protect\hypertarget{blinding-placebo}{}
\textbf{Blinding}: Keeping subjects from knowing which group they are
in: give a ``fake'' treatment (\textbf{placebo}) to the control group.

\begin{itemize}
\tightlist
\item
  Knowing if you are being treated or not might be confounding.
\item
  Example: Pain killer
\end{itemize}

\begin{center}
\begin{tabular}{c|cc}
&Study 1 & Study 2  \\[-3pt]
&(w.~placebo) & (w/o.~placebo)\\
\hline
%& Rate &Rate \\
Treatment & 30\%  & 31\% \\
Control   & 31\%  & 20\% \\
\end{tabular}
\end{center}

Is the drug effective? Which study should we look at?
\end{frame}

\begin{frame}{Single Blind v.s. Double Blind}
\protect\hypertarget{single-blind-v.s.-double-blind}{}
\textbf{Single Blind}: Only the subject are kept from knowing which
treatment they received.

\textbf{Double Blind}: Neither the subject nor those who evaluate the
outcome (e.g., diagnosing the disease) know which treatment the subject
receive

\begin{itemize}
\tightlist
\item
  To avoid the bias in the evaluation, especially in the borderline
  cases.
\end{itemize}
\end{frame}

\begin{frame}{Example: Polio Vaccine Results}
\protect\hypertarget{example-polio-vaccine-results}{}
\begin{center}
\underline{Randomized controlled double-blind experiment}\\
\begin{tabular}{ccc}
  & Size  & Rate {\scriptsize(per 100,000)} \\
\makebox[2cm]{Treatment} & \makebox[2cm]{200,000} & \makebox[2cm]{28} \\
\makebox[2cm]{Control} & \makebox[2cm]{200,000} & \makebox[2cm]{71} \\
\makebox[2cm]{No consent} & \makebox[2cm]{350,000} & \makebox[2cm]{46}
\end{tabular}
\end{center}

\begin{enumerate}
\item
  Which pair of rates shows the vaccine is effective?

  \begin{center} (a) 28 and 71 \qquad (b) 28 and 46 \qquad (c) 71 and 46\end{center}
\item
  Neither the control group nor the non-consent group got the vaccine,
  but the no-consent group had a lower rate of polio. Why?
\item
  In a badly designed test where all volunteers are assigned to the
  treatment group, and all non-volunteers are to the control group,
  which pair of rates are the most likely rates for the two groups?

  \begin{center} (a) 28 and 71 \qquad (b) 28 and 46 \qquad (c) 71 and 46\end{center}
\end{enumerate}
\end{frame}

\begin{frame}{Comparison of the Two Polio Vaccine Experiments}
\protect\hypertarget{comparison-of-the-two-polio-vaccine-experiments}{}
\begin{center}
\tabcolsep=3pt
\begin{tabular}{lcc||ccc}
\multicolumn{3}{c||}{Randomized controlled} &\\
\multicolumn{3}{c||}{double-blind experiment} &
\multicolumn{3}{c}{NFIP study} \\\hline
 & Size & Rate &  Grade & Size & Rate \\
Treatment  & 200,000 & 28 &  2 (vaccine) & 225,000 & 25 \\
Control    & 200,000 & 71 &  1 \& 3 (control) & 725,000 & 54 \\
No consent & 350,000 & 46 &  2 (no consent) & 125,000 & 44
\end{tabular}
\end{center}

\begin{enumerate}
\setcounter{enumi}{3}
\item
  The two vaccine groups had similar rates of polio, 28 and 25. \newline
  The two no-consent groups also had similar rates of polio, 46 and
  44.\newline Why are the rates of polio in the two control groups so
  different?
\item
  In the NFIP study, neither the control group nor the no-consent group
  got the vaccine. Yet the no-consent group had a lower rate of polio.
  Why?\medskip
\item
  Can one conclude that the vaccine is effective by comparing the 44 in
  the NFIP study with the 25 in the vaccine group?

  \begin{itemize}
  \item
    \uncover<2->{\alert{No! volunteers/non-volunteers comparison}}
  \end{itemize}
\end{enumerate}
\end{frame}

\begin{frame}{Historical and Contemporaneous Controls}
\protect\hypertarget{historical-and-contemporaneous-controls}{}
\begin{itemize}
\tightlist
\item
  Sometimes randomized controlled experiments are difficult, expensive,
  or unethical to do.
\item
  One commonly used (but poor) alternative is to compare patients using
  new treatment with \textbf{historical controls}: patients treated in
  the old way \textbf{in the past}.
\item
  In contrast, in randomized controlled experiments, subjects in both
  groups are chosen from a population at the same time period
  (\textbf{contemporaneous controls}).
\item
  Historical controls may not be reliable because the treatment group
  and the historical control group may differ in important ways.
\end{itemize}
\end{frame}

\end{document}
